\section{Malware and Carrier Taxonomy}\label{sec:malware-taxonomy}

Malware terminology becomes unreliable when a delivery artifact, a propagation
mechanism, and a payload objective are treated as interchangeable labels. For
this study, an attachment is described on five orthogonal axes. The first is
the \emph{replication mechanism}, such as infection of a host program or
self-propagation between systems after execution. The second is the \emph{initial-access
mechanism}, such as a spearphishing attachment, a deceptive download, or a file
placed in a shared location. The third is the \emph{operational role}, such as
downloader, dropper, or stager. The fourth is the \emph{payload objective}, such
as encryption of user data, surveillance, or credential theft. The fifth is
the \emph{carrier property}, meaning a feature of the submitted name, declared
type, byte grammar, resource demand, metadata, or reconstruction result that
\AFD{} can actually observe.

Only the fifth axis directly determines the file-defense verdict. The first
four axes define threat context and corpus strata. This rule prevents a
blocked Portable Executable file named as an image from being reported as a
detected Trojan unless independent ground truth establishes that behavioral
label. It also prevents an accepted image from being reported as harmless. An
accepted image satisfies a selected output profile and policy. It may still
depict harmful instructions, carry adversarial meaning, or preserve information
encoded in pixels.

\subsection{Replication mechanism}

A file-infector virus modifies another program so that execution of the host
also propagates a version of the virus~\cite{cohen1987}. The observable carrier
may therefore be an executable prefix, an altered executable host grammar, an
appended body, or a composite object that remains executable to one consumer.
\AFD{} can reject such an object when it is an unsupported executable, when name
and bytes disagree, or when complete incompatible parses violate the
true-polyglot ambiguity policy. The pipeline does not emulate the candidate and
cannot establish that its code would infect another program.

A worm is code that self-propagates to another machine or account after
execution, without requiring a user to initiate each subsequent copy. It may
use network reachability, service vulnerabilities, or credentials. A user may
still initiate the first execution. Spafford's analysis shows why this
propagation model is not reducible to file infection~\cite{spafford1989}. An
attachment may be the initial entry point, but later propagation may use no
attachment. Consequently, the measurable event is the blocking of an active or
unsupported carrier, not the classification of worm replication.

\subsection{Initial access and represented function}

Trojan terminology concerns represented purpose and actual purpose. In an
attachment workflow, the relevant observable form is often masquerading. A
file may use a benign-looking icon or name, a double extension, a supplied
media type inconsistent with its bytes, or a bidirectional control character
that changes visual ordering. MITRE ATT\&CK records malicious-file execution as
a user-execution technique. It separately records right-to-left override,
double file extensions, and masqueraded file types as masquerading
sub-techniques~\cite{mitreT1204002,mitreT1036002,mitreT1036007,
mitreT1036008}. These observations justify name normalization and evidence
convergence. They do not prove that every disagreement is malicious. The
fail-closed release policy blocks ambiguity because it cannot establish a
supported media identity, not because it has inferred intent.

Spearphishing attachment is an initial-access mechanism rather than a malware
family. ATT\&CK describes it as targeted electronic delivery that commonly
relies on user execution~\cite{mitreT1566001}. The same mechanism may deliver a
Trojan, downloader, dropper, stager, spyware implant, or ransomware program.
Conversely, the same payload can arrive through a website, removable media, or
an already compromised account. Corpus records therefore store delivery context
separately from the observed attachment property.

\subsection{Operational roles}

This study distinguishes three staged-delivery roles. A \emph{downloader}
obtains a follow-on payload from another source after it executes. A
\emph{dropper} carries an embedded payload and writes, installs, or otherwise
places that payload after execution. A \emph{stager} is a first component whose
role is to prepare an execution context or transfer control to a follow-on
stage. A stager can implement downloading or dropping, but the terms are not
synonyms. Downloader and dropper describe how the follow-on material is
obtained. Stager describes the component's position in a multi-stage chain.

Scripts, shortcuts, executables, active documents, and packages can carry any
of these roles. They are unsupported by the media-only policy and can be
blocked before the role is observed. An ordinary image containing a visible URL
is not thereby a downloader. Likewise, an exploit that causes a vulnerable
decoder to initiate network activity remains an exploit-bearing media case
unless independent execution evidence establishes a downloader or stager role.

\subsection{Payload objective}

Ransomware has the objective of denying access to data or systems, commonly
with a demand for payment. Research systems that claim ransomware detection
observe file-system and desktop behavior during execution~\cite{kharraz2016,
scaife2016}. \AFD{} observes neither mass file transformation nor ransom-note
presentation. A corpus sample with independently established ransomware ground
truth may be used to ask whether its carrier is blocked. The result must be
reported as carrier containment. It cannot be generalized to ransomware that
arrives through an allowed representation, exploits a retained codec stream,
or is already resident on the endpoint.

Spyware has a collection or disclosure objective. An executable spyware
carrier is blocked for the same media-policy reasons as another executable. A
passive metadata field does not execute collection or disclosure behavior and
is therefore not spyware. The spyware label requires independent evidence of a
surveillance program or its behavior, not merely the presence of location,
author, device, or application data in a media file.

\subsection{Privacy leakage}

Passive metadata can leak privacy when a user shares the file. The pipeline can
remove named fields such as source-application records or image location data
when the output profile excludes them. This is metadata exclusion and privacy
leakage reduction, not spyware detection. It does not establish that a program
would record input, capture a screen, or transmit stored data. It also does not
remove information visibly or covertly encoded into retained pixel or sample
values.

\subsection{Carrier properties}

Carrier properties translate the threat taxonomy into testable predicates.
An \emph{unsupported active class} is an executable, script, archive, active
document, shortcut, package, or other grammar outside the frozen media-only
allowlist. A \emph{classifier disagreement} occurs when normalized name,
supplied media type, and bounded byte evidence do not converge on a compatible
concrete profile. A \emph{malformed supported object} claims an allowed media
family but violates lengths, offsets, checksums, mandatory elements, nesting,
or terminal boundaries. A \emph{resource-exhaustion carrier} exceeds a bound
on input or output size, decoded dimensions, frame count, duration, component
count, nesting, expansion, memory, cancellation, or time.

An \emph{exploit-bearing media object} is a supported-looking carrier with a
reproducible structure or codec payload that triggers an identified vulnerable
path. Malformation and exploit are not synonyms. Some malformed files are
benign corruption, and some exploits are valid according to a format
specification. The permitted evaluation claim names the frozen parser or
decoder path and the sample outcome. Rejecting one regression does not establish
coverage of all vulnerabilities in the codec family. Invoking a decoder under
limits also does not make that decoder intrinsically safe.

A \emph{true polyglot} is a whole octet string that is completely accepted,
from offset zero through the same terminal boundary, by at least two
incompatible parsers. This definition requires complete acceptance of the
entire file under both grammars. A valid media object followed by bytes that
form or resemble a second object is an \emph{appended or trailing-object
carrier}, not a true dual-format polyglot, unless both parsers accept the whole
string completely. Prefix camouflage likewise places unsupported bytes before
a later media signature without by itself establishing dual-format validity.
An overlapping construction is a true polyglot only when the complete
acceptance condition holds. A nested active object occupies a format-defined
location inside a container and is a separate carrier property.

The defense mechanism depends on the class. Parsing must start at zero to
reject prefix camouflage and terminate at the exact file boundary to reject
unconsumed trailing bytes. A true polyglot requires reports from at least two
incompatible complete parsers. Location-aware component allowlists govern
nested objects. Substring scanning inside compressed payloads is not sufficient
evidence of a nested object and can create false attributions.

Structural and semantic reconstruction support different conclusions. A
structural path can rebuild container structure, remove excluded metadata, and
check complete boundaries while retaining encoded frames or samples. It can
therefore support container-smuggling and metadata claims. It cannot claim to
neutralize a secondary interpretation inside retained codec bytes. A semantic
path decodes bounded pixels, frames, or samples and encodes a new object without
copying source container or compressed payload bytes. This supports a stronger
sample-specific neutralization claim, conditional on decoder and encoder
behavior and on independent validation of the result. It still does not prove
that the semantic content is harmless.

\begin{landscape}
\begingroup
\footnotesize
\setlength{\tabcolsep}{3pt}
\begin{longtable}{L{2.3cm}L{2.2cm}L{3.8cm}L{4.4cm}L{4.4cm}}
\caption{Orthogonal malware and carrier taxonomy used for corpus annotation
and claim control. A sample may occupy several rows because operational role
and carrier form are independent.}
\label{tab:malware-carrier-taxonomy}\\
\toprule
Term & Taxonomic axis & Observable attachment property & Permitted claim & Excluded inference \\
\midrule
\endfirsthead
\toprule
Term & Taxonomic axis & Observable attachment property & Permitted claim & Excluded inference \\
\midrule
\endhead
\midrule
\multicolumn{5}{r}{Continued on next page}\\
\endfoot
\bottomrule
\endlastfoot
Virus & Replication & Executable host, executable prefix, appended body, or incompatible grammar & Type, grammar, or carrier policy blocked the object & The bytes would infect another program \\
Worm & Replication & Unsupported executable, script, shortcut, package, or active document & Unsupported active carrier was blocked & The sample would self-propagate after execution \\
Trojan & Represented function & Name, media type, and bytes disagree, or active content is presented as media & Masquerading evidence or unsupported grammar was blocked & The file's represented and actual purposes were behaviorally established \\
Downloader & Operational role & Executable, script, shortcut, package, or active document & Media-only policy blocked the carrier before execution & The component would obtain a remote follow-on payload \\
Dropper & Operational role & Executable, script, shortcut, package, installer, or active document & Media-only policy blocked the carrier before execution & The component contained and would place an embedded payload \\
Stager & Operational role & Executable, script, shortcut, package, installer, or active document & Media-only policy blocked the carrier before execution & The component would prepare or transfer control to a later stage \\
Ransomware & Payload objective & Independently labelled sample carried by executable or script & The labelled carrier was blocked & Ransomware behavior was classified \\
Spyware & Payload objective & Independently labelled executable or active carrier & The labelled carrier was blocked under the media-only policy & Surveillance behavior was observed or classified \\
Passive metadata & Privacy leakage & Named location, author, device, application, comment, or thumbnail field & The specified field was absent from the parsed output & The file was spyware or fully de-identified \\
Exploit-bearing media & Carrier effect & Reproducible malformed structure or decoder-triggering sample & Frozen path blocked or reconstructed this sample & All exploits for the format were neutralized \\
Masquerading & Initial access or evasion & Double extension, bidi control, false media type, or byte disagreement & Name policy or type convergence blocked ambiguity & Malicious intent was proven \\
Prefix camouflage & Carrier structure & Unsupported bytes precede a later media signature & Zero-offset parsing or type convergence blocked the object & The whole file was valid under two formats \\
Appended or trailing object & Carrier structure & Bytes remain after the supported object's exact terminal boundary & Complete-consumption policy blocked or complete reconstruction removed the suffix & The whole file was a true polyglot \\
True polyglot & Carrier structure & The whole octet string is completely accepted by at least two incompatible parsers & Ambiguity policy blocked the object, or a qualified semantic path emitted a uniquely accepted output & Every secondary meaning in retained semantic data disappeared \\
Nested active object & Carrier structure & A parsed active object occurs at a format-defined location excluded by policy & Location-aware policy removed the object through reconstruction or blocked the carrier & A substring match proved an embedded object \\
Stegano\-graphic content & Retained semantics & Information encoded in pixels, frames, or samples & Outside the structural and behavioral claim & Accepted bytes are harmless \\
\end{longtable}
\endgroup
\end{landscape}

\subsection{Corpus annotation and reporting rules}

Each corpus record should identify provenance, ground-truth basis, concrete
format profiles, carrier class, expected action, and expected reason family.
Behavioral labels require independent evidence such as a reproducible analysis
or expert adjudication. They must not be inherited from a filename, an upload
directory, or a scanner label alone. Carrier labels require parser reports or a
reproducible mutation recipe. A sample can therefore be a ransomware-labelled
executable, delivered by spearphishing, masquerading as a JPEG, with an
appended-media signature. Those are four annotations, not four synonyms.

Aggregate reporting follows the same discipline. Results are stratified by
carrier property and expected reason family. A blocked ransomware-labelled
executable contributes to unsupported-active-carrier containment. A malformed
JPEG contributes to structural rejection. A semantically rebuilt image-based
true polyglot contributes to polyglot neutralization only if complete parser
reports establish whole-string acceptance under at least two incompatible
grammars and every reconstruction prerequisite is recorded. A removed trailing
object contributes instead to complete-consumption or suffix-removal evidence.
No aggregate ``malware detection rate'' is computed from these heterogeneous
outcomes. Empirical counts and rates remain unset until the frozen corpus and
controlled server run satisfy the pre-specified evidence gates.
